\chapter{Menyelesaikan Soal Cerita}\label{ch:g3:problems}

Menguasai empat operasi itu satu hal; mengetahui \emph{yang mana}
yang diperlukan sebuah cerita adalah hal lain. Bab ini memberi cara
untuk mengerjakan soal cerita --- empat langkah yang sama setiap
kali --- lalu mengajarkan cara membaca keterangan dari tabel dan
piktogram.

\section{Empat langkah}

\begin{method}[Menyelesaikan soal cerita]\label{met:g3:problems:steps}
\begin{enumerate}
  \item \emph{Bacalah} soalnya dua kali; katakan dengan kata-katamu
        sendiri apa yang ditanyakan dan apa yang diketahui.
  \item \emph{Pilihlah} operasinya: buatlah sketsa cepat atau gambar
        batang kalau pilihannya belum jelas.
  \item \emph{Hitunglah}, secara bersusun atau di kepala, dengan
        taksiran sebagai jaring pengaman.
  \item \emph{Jawablah} dengan kalimat lengkap beserta satuannya ---
        lalu periksa bahwa jawabannya masuk akal (berat seorang anak
        bukan $300$ kg).
\end{enumerate}
\end{method}

\begin{example}[Operasi yang mana?]\label{ex:g3:problems:which}
Kata-katanya memberi petunjuk, sketsanya yang memutuskan:
\begin{itemize}
  \item \emph{digabungkan / seluruhnya}: penjumlahan;
  \item \emph{diambil / berapa lebihnya / apa yang hilang}:
        pengurangan;
  \item \emph{sekian kelompok berisi sekian / sekian kali}:
        perkalian;
  \item \emph{dibagi rata / berapa kelompok}: pembagian.
\end{itemize}
\end{example}

\begin{omfigure}
\begin{tikzpicture}[scale=0.62]
  % batang penjumlahan
  \draw[thick, fill=omDef!20] (0,0) rectangle (3.4,0.8);
  \draw[thick, fill=omThm!20] (3.4,0) rectangle (5.8,0.8);
  \node[font=\small] at (1.7,0.4) {$26$};
  \node[font=\small] at (4.6,0.4) {$17$};
  \draw[Stealth-Stealth] (0,-0.35) -- (5.8,-0.35);
  \node[below, font=\small] at (2.9,-0.4) {? seluruhnya: $26 + 17$};
  % batang pembanding
  \begin{scope}[xshift=9.4cm]
  \draw[thick, fill=omDef!20] (0,0.9) rectangle (5.2,1.7);
  \node[font=\small] at (2.6,1.3) {$52$};
  \draw[thick, fill=omThm!20] (0,0) rectangle (3.5,0.8);
  \node[font=\small] at (1.75,0.4) {$35$};
  \draw[Stealth-Stealth] (3.5,0.4) -- (5.2,0.4);
  \node[below, font=\small] at (2.6,-0.15)
    {? lebihnya: $52 - 35$};
  \end{scope}
\end{tikzpicture}

{\small Gambar batang membuat operasinya terlihat: bagian yang
berdampingan meminta penjumlahan; dua batang yang dibandingkan
meminta pengurangan.}
\end{omfigure}

\begin{example}[Soal dua langkah]\label{ex:g3:problems:twostep}
``Maya membeli $3$ bungkus berisi $6$ botol lalu meminum $2$ botol.
Berapa botol yang tersisa?''
\begin{enumerate}
  \item Yang ditanyakan: banyaknya botol yang tersisa. Yang
        diketahui: $3$ bungkus berisi $6$; $2$ botol diminum.
  \item Sketsa: $3$ kelompok berisi $6$, lalu ambil $2$. Dua operasi!
  \item Hitung: $3 \times 6 = 18$, lalu $18 - 2 = 16$.
  \item Kalimat: tersisa $16$ botol.
\end{enumerate}
Jangan pernah berhenti setelah operasi pertama: bacalah lagi
pertanyaannya untuk melihat apakah ceritanya sudah selesai.
\end{example}

\begin{example}[Bilangan yang tidak berguna]\label{ex:g3:problems:trap}
``Seorang pembuat roti, yang berumur $52$ tahun, memanggang $120$ roti
lalu menjual $85$. Berapa roti yang tersisa?'' Umurnya ada di sana
untuk mengecohmu: jawabannya hanya memerlukan $120 - 85 = 35$ roti.
Sebelum menghitung, pilahlah selalu bilangannya: mana yang
sungguh-sungguh dipakai pertanyaannya?
\end{example}

\section{Membaca tabel dan piktogram}

\begin{example}[Sebuah tabel]\label{ex:g3:problems:table}
Buku yang dipinjam di perpustakaan:
\begin{center}
\renewcommand{\arraystretch}{1.3}
\begin{tabular}{l|cccc}
 & komik & novel & sains & \omterm{def:g1:addition:def}{jumlah} \\
\hline
minggu 1 & $14$ & $9$ & $6$ & $29$ \\
minggu 2 & $11$ & $12$ & $8$ & $31$ \\
\end{tabular}
\end{center}
Membaca: minggu 2, novel: $12$. Menghitung dari tabel: dalam dua
minggu itu, $14 + 11 = 25$ komik dipinjam; minggu yang paling sibuk
adalah minggu 2 ($31 > 29$).
\end{example}

\begin{example}[Sebuah piktogram]\label{ex:g3:problems:pictogram}
Dalam piktogram, satu lambang mewakili \omterm{def:g1:addition:def}{jumlah} yang tetap --- bacalah
kuncinya dulu!

\begin{omfigure}
\begin{tikzpicture}[scale=0.62]
  \node[left, font=\small] at (0,2.4) {Senin};
  \foreach \i in {0,1,2} \fill[omDef] (0.5+\i*0.8,2.4) circle (0.26);
  \node[left, font=\small] at (0,1.2) {Selasa};
  \foreach \i in {0,1,2,3,4} \fill[omDef] (0.5+\i*0.8,1.2) circle (0.26);
  \node[left, font=\small] at (0,0) {Rabu};
  \foreach \i in {0,1} \fill[omDef] (0.5+\i*0.8,0) circle (0.26);
  \fill[omDef] (0.5+1.6,0) circle (0.26);
  \node[right, font=\small] at (5.4,1.2)
    {kunci: satu titik $= 4$ kue terjual};
\end{tikzpicture}

{\small Sebuah piktogram: Senin menunjukkan $3$ titik, jadi
$3 \times 4 = 12$ kue; Selasa $5$ titik, $20$ kue; Rabu $3$ titik,
$12$ kue.}
\end{omfigure}
\end{example}

\section{Latihan}

\begin{exercise}[$\star$]\label{exo:g3:problems:1}
Untuk setiap cerita, cukup \emph{sebutkan} operasinya (jangan
dihitung): sekawanan $48$ domba bergabung dengan kawanan $37$ domba;
\ $56$ ceri dibagi rata kepada $7$ anak; \ $9$ kotak berisi $8$
pensil; \ Ali punya $71$ kartu lalu memberikan $18$.
\end{exercise}

\begin{exercise}[$\star$]\label{exo:g3:problems:2}
Selesaikan dengan empat langkah: ``Sebuah tempat parkir berisi $246$
mobil pada pagi hari; $158$ mobil lagi datang siang hari. Berapa mobil yang ada sekarang?''
\end{exercise}

\begin{exercise}[$\star$]\label{exo:g3:problems:3}
Selesaikan: ``Seutas tali sepanjang $90$ m dipotong menjadi potongan
$9$ m. Berapa banyak potongannya?''
\end{exercise}

\begin{exercise}[$\star$]\label{exo:g3:problems:4}
Selesaikan: ``Harga satu karcis $8$. Berapa harga $26$ karcis?''
\end{exercise}

\begin{exercise}[$\star$]\label{exo:g3:problems:5}
Gambarlah gambar batangnya, lalu selesaikan: ``Emi punya $64$ stiker,
stiker Bayu $29$ lebih sedikit. Berapa stiker Bayu?''
\end{exercise}

\begin{exercise}[$\star$]\label{exo:g3:problems:6}
Dua langkah: ``Satu peti berisi $45$ apel. Seorang pedagang menerima
$6$ peti lalu menjual $178$ apel. Berapa apel yang tersisa?''
\end{exercise}

\begin{exercise}[$\star$]\label{exo:g3:problems:7}
Dengan tabel pada \cref{ex:g3:problems:table}: berapa buku sains
yang dipinjam dalam dua minggu itu? Berapa novel lebih banyak pada
minggu 2 daripada pada minggu 1?
\end{exercise}

\begin{exercise}[$\star$]\label{exo:g3:problems:8}
Dengan piktogram pada \cref{ex:g3:problems:pictogram}: berapa kue
yang terjual dalam tiga hari itu? Pada hari apa kue terjual paling
banyak?
\end{exercise}

\begin{exercise}[$\star$]\label{exo:g3:problems:9}
Soal ini punya satu bilangan yang tidak berguna --- carilah, lalu
selesaikan: ``Sebuah buku setebal $210$ halaman berharga $12$. Leo
membaca $35$ halaman setiap hari. Berapa hari yang ia perlukan untuk menyelesaikan buku itu?''
\end{exercise}

\begin{exercise}[$\star\star$]\label{exo:g3:problems:10}
``Sebuah keluarga yang terdiri atas $2$ orang dewasa dan $3$ anak
mengunjungi museum. Karcis dewasa berharga $9$, karcis anak $4$.
Mereka membayar dengan uang $50$. Berapa kembaliannya?'' (Tiga
langkah: dewasa, anak, lalu kembaliannya.)
\end{exercise}

\begin{exercise}[$\star\star$]\label{exo:g3:problems:11}
Karanglah sendiri soal dua langkah yang jawabannya $100$, tulislah
penyelesaiannya dengan empat langkah, lalu ujilah pada teman sekelas.

\end{exercise}
